\chapter{Research Methodology}\label{ch:methodology}

The four core chapters of this thesis study different problems: 3D occupancy perception, language-based reasoning, test-time adaptation of driving policies, and robust robot manipulation. This chapter describes the research framework behind them. The framework is a reflective reconstruction: it was not designed in advance and then applied, but distilled after the fact from the author's projects, including co-authored projects in which the author participated without leading them. It is presented here for two reasons. First, it makes explicit, and justifies, the methodological choices of Chapters~\ref{ch:occnet}--\ref{ch:kai0}, so that those chapters can concentrate on technical content. Second, it offers a framework that other researchers working on physical agents may adapt, subject to the limits discussed in Sections~\ref{sec:method:cycle} and~\ref{sec:method:validity}.

The chapter describes the research paradigm (Section~\ref{sec:method:paradigm}), the conceptual framework through which the ``curse of reality'' is analysed (Section~\ref{sec:method:framework}), the research cycle and the lineage of the author's publications (Sections~\ref{sec:method:cycle} and~\ref{sec:method:lineage}), the design principles distilled from them (Section~\ref{sec:method:principles}), the evaluation methodology (Section~\ref{sec:method:evaluation}), a decision procedure that combines these elements (Section~\ref{sec:method:decision}), and the practices adopted for reproducibility, ethics and validity (Sections~\ref{sec:method:practice} and~\ref{sec:method:validity}).

% ==============================================================================
\section{Research Paradigm}\label{sec:method:paradigm}

Research on physical agents is, in large part, a design science~\cite{simon1996sciences,hevner2004design}. The objects of study are artefacts (representations, models, datasets and evaluation protocols), and knowledge is produced by building such artefacts and testing them against the environment they are meant to operate in. A claim such as ``3D occupancy is a better scene representation than 3D boxes for planning'' cannot be settled by argument alone; it is settled by constructing both representations, training systems on them, and measuring how the systems behave on data that reflect the real world. The research in this thesis adopts that paradigm and, within it, two commitments.

\paragraph{Commitment 1: benchmarks accelerate the research cycle.} A research question becomes tractable for a community only when it is measurable: when there is a task definition, a public dataset and a metric that many groups can use. The construction of benchmarks, and of public challenges around them, is therefore treated as a research contribution in its own right, not as a preliminary to the ``real'' work. More than half of the eighteen publications summarised in Table~\ref{tab:method:papers} introduce a dataset, a benchmark split or a data platform.

\paragraph{Commitment 2: a good module must scale to real-world use.} A method that works only at a small scale, only with privileged information, or only at a computational cost that no deployed system can afford does not solve the problem it claims to solve. Methods are therefore judged on whether they can be trained on large data, whether they use inputs available at deployment, and whether their cost fits the deployment budget. This commitment appears, for example, in the preference for camera-based perception that uses LiDAR at training time only, studied in co-authored work~\cite{huang2023vcd}; in the single test-time gradient step of Chapter~\ref{ch:centaur}; and in the post-training resource constraints studied in Chapter~\ref{ch:kai0}.

% ==============================================================================
\section{Conceptual Framework: the Curse of Reality}\label{sec:method:framework}

\subsection{Formulation}\label{sec:method:formulation}

Let $s$ denote the state of the world (the scene and the agents in it), $o$ the sensor observation, and $a$ the action. A simplified single-time-step description of imitation-learning data is
\begin{equation}
  p_{\mathrm{train}}(o,a) \;=\; \sum_{s} p_{\mathrm{train}}(s)\; p_{\mathrm{train}}(o \mid s)\; \pi^{*}(a \mid o),
  \label{eq:method:train}
\end{equation}
where $\pi^{*}$ is the demonstrator, here assumed to condition on $o$, and the sum is replaced by an integral for continuous states. This equation describes observation--action pairs, not a distribution over closed-loop trajectories. A learned policy $\pi$ is deployed in a world with its own state distribution $p_{\mathrm{deploy}}(s)$ and sensing model $p_{\mathrm{deploy}}(o \mid s)$. In closed loop, the states it visits follow its own state distribution $d_{\pi}$ rather than the demonstrator's $d_{\pi^{*}}$, and its actions are executed by hardware as $a_{\mathrm{exec}} \sim p(a_{\mathrm{exec}} \mid a)$. The curse of reality concerns differences between training and deployment that finite training data may not cover, and that rare states, individually improbable but collectively frequent, are among the most consequential~\cite{liu2024curse}: the long tail is never exhausted, and in closed loop small errors compound~\cite{ross2011dagger}.

\subsection{Four kinds of shift}\label{sec:method:shifts}

The author's work suggests four diagnostic categories (Fig.~\ref{fig:method:shift}). They cover sensing and representation limits, unfamiliar situations, policy-induced state visitation, and differences between training and execution. The categories are not mutually exclusive and do not correspond one-to-one to factors of Eq.~\eqref{eq:method:train}. They help organise failures and possible responses, rather than define a partition or a sequence of processing stages.

\begin{figure}[t]
  \centering
  \tikzfig{ch-methodology/figures/shift_taxonomy_revised}
  \caption[Four diagnostic categories of shift]{Four non-exclusive diagnostic categories of shift. The rows associate symptoms with possible methodological responses and the chapters that study them; they are not a probabilistic factorisation or an execution pipeline.}
  \label{fig:method:shift}
\end{figure}

\paragraph{Perceptual shift: what is out there?} The sensing model $p(o \mid s)$ changes, or the state contains entities that the representation cannot express. Roads are not flat: in the OpenLane benchmark, only about 20\% of frames vary by less than \SI{0.5}{\metre} in height, and more than half vary by over \SI{1}{\metre}~\cite{chen2022persformer}. Irregularly shaped vehicles such as trucks, trailers and construction vehicles are poorly described by bounding boxes; supervising them with occupancy instead of voxelised boxes improves their segmentation by 2.21 mIoU at \SI{0.5}{\metre} resolution and by 3.51 mIoU at \SI{0.25}{\metre}~\cite{sima2023occnet}. Sensing conditions also change: a camera-based detector loses much of its accuracy from day to night~\cite{li2023bevdevils}, a fusion detector must survive the partial failure of a sensor~\cite{gao2023sdf}, and a driving policy may be deployed on a vehicle whose camera configuration differs from that of its training data, as in the zero-shot transfer from nuScenes to Waymo studied in Chapter~\ref{ch:drivelm}.

\paragraph{Semantic shift: what does it mean?} The state distribution $p(s)$ changes: situations or compositions of agents that are rare or absent in training call for knowledge that the training labels do not contain. In a CARLA split containing pedestrians that were absent from training, the behaviour accuracy of DriveLM-Agent dropped from 59.63 to 4.59 (and that of TransFuser++ from 70.19 to 8.72); adding one explicit question about pedestrians to the reasoning graph raised it to 27.04 (Chapter~\ref{ch:drivelm}). Measuring such generalisation is itself difficult: the follow-up audit DriveBench showed that vision-language models can keep most of their score even when the image is removed, answering from textual priors rather than from the scene~\cite{xie2025vlms}.

\paragraph{Behavioural shift: what will my action cause?} In closed loop, the policy's own state distribution $d_{\pi}$ departs from the demonstrator's, and rare, safety-critical states arise that demonstrations seldom contain. Chapter~\ref{ch:centaur} evaluates uncertainty-based test-time training in non-reactive NAVSIM. Its fixed-bank fallback and UniAD optimisation records suggest progress--safety trade-offs, but the fallback aggregation requires verification. The model-selected \texttt{navsafe} split gives category-level stress tests rather than an unbiased aggregate performance estimate.

\paragraph{Deployment shift: does training match use?} The way a policy is trained differs from the way it is used. Chapter~\ref{ch:kai0} studies this in robot manipulation through human demonstrations, the learned policy and trajectories executed on hardware. It uses coverage deficiency, temporal mismatch and failure cascade to describe possible failure mechanisms. This is a diagnostic account, not a measured decomposition of performance loss among three distributions. The authors also report, as an observation rather than a controlled result, that success rates varied between 20\% and 60\% with data quality under otherwise identical settings~\cite{sima2026kai0}.

\subsection{Methodological levers}

Each category suggests possible responses, called here methodological levers. Chapter~\ref{ch:occnet} studies richer geometric representations; Chapter~\ref{ch:drivelm} exposes intermediate reasoning for supervision and inspection; Chapter~\ref{ch:centaur} studies test-time policy updates without new expert labels; and Chapter~\ref{ch:kai0} studies training and execution alignment under resource constraints. These are examples rather than exclusive assignments of one remedy to each shift. Two further levers cut across all four. The first is data: scaling or curating the training data is always available and was studied in the co-authored AgiBot World~\cite{bu2025agibot_iros}, although it is not the subject of a core chapter. The second is measurement: none of the other levers can be pulled, or its effect verified, without a benchmark that exposes the corresponding shift.

% ==============================================================================
\section{The Research Cycle}\label{sec:method:cycle}

Looking back over the author's publications, the research can be described as a recurring cycle of six stages (Fig.~\ref{fig:method:cycle}). The cycle is a description of how the work proceeded, reconstructed in hindsight; it is suggested, not prescribed, as a framework for similar research, and its limits are discussed at the end of this section.

\begin{figure}[t]
  \centering
  \tikzfig{ch-methodology/figures/research_cycle_revised}
  \caption[The benchmark-driven research cycle]{The benchmark-driven research cycle reconstructed from the author's publications. Grey annotations give concrete artefacts at each stage; works marked $^{*}$ are co-authored rather than led by the author. The cycle closes when residual failures motivate the next round of formalisation.}
  \label{fig:method:cycle}
\end{figure}

\paragraph{Stage 1: observe.} A project starts from a failure observed, or anticipated, in the field rather than from a gap in a leaderboard. Examples are obstacles that bounding boxes cannot describe (Chapter~\ref{ch:occnet}), driving decisions that cannot be explained or that break under a change of sensors or scenes (Chapter~\ref{ch:drivelm}), rare safety-critical scenes (Chapter~\ref{ch:centaur}), and robot-policy failures associated with inconsistent demonstrations or execution (Chapter~\ref{ch:kai0}).

\paragraph{Stage 2: formalise.} The failure is turned into a task with a representation and a metric: semantic occupancy with voxel-level mIoU and flow (Chapter~\ref{ch:occnet}); graph visual question answering with language, behaviour and motion metrics (Chapter~\ref{ch:drivelm}); the PDM score of NAVSIM together with a safety-critical split (Chapter~\ref{ch:centaur}); and real-world task success and recovery (Chapter~\ref{ch:kai0}). The formalisation fixes what counts as progress, so it deserves as much care as the method.

\paragraph{Stage 3: benchmark.} Data are collected or re-annotated at a scale that makes the task meaningful, usually by an automatic labelling pipeline followed by human verification. Benchmarks built in this way include OpenLane (200K frames, over 880K 3D lanes)~\cite{chen2022persformer}; OpenLane-V2 (2,000 scenes, 2.1M instances and 1.9M relations for road topology)~\cite{wang2023openlanev2}; OpenOcc (34,149 frames, over 1.4 billion annotated voxels)~\cite{sima2023occnet}; OpenScene (more than 120 hours of driving with occupancy labels)~\cite{sima2023openscene}; DriveLM-nuScenes (4,871 frames with 91.4 question--answer pairs per frame) and DriveLM-CARLA (about 1.6M question--answer pairs)~\cite{sima2024drivelm}; the \texttt{navsafe} split (229 safety-critical frames in 10 categories, Chapter~\ref{ch:centaur}); and AgiBot World (1,001,552 trajectories, 2,976.4 hours, 217 tasks)~\cite{bu2025agibot_iros}.

\paragraph{Stage 4: baseline.} A simple, unified and scalable baseline is built on the benchmark. It shows that the task is solvable and gives the community a reference point. Deliberately simple baselines are one choice, for example the BLIP-2-based DriveLM-Agent of Chapter~\ref{ch:drivelm}; strong unified ones are another, for example the co-authored BEVFormer, which reached 56.9 NDS on the nuScenes test set, 9.0 points above the previous state of the art~\cite{li2022bevformer}, and UniAD, which chained all driving sub-tasks towards planning~\cite{hu2023uniad}.

\paragraph{Stage 5: community.} Code, data and models are released, and public challenges are organised on top of the benchmark. The author co-organised the 3D Occupancy Prediction challenge at CVPR 2023 and the ``Driving with Language'' track of the Autonomous Grand Challenge at CVPR 2024, which attracted 152 teams from 14 countries and regions with 978 submissions.\footnote{Challenge pages: \url{https://opendrivelab.com/challenge2023/} and \url{https://opendrivelab.com/challenge2024/}.} OpenScene provided the data for the end-to-end driving and predictive world-model tracks at CVPR 2024 and for the NAVSIM track at CVPR 2025, and AgiBot World hosted its own challenge in 2025. A challenge is not an end in itself: it stress-tests the benchmark design, exposes shortcuts, and produces strong baselines faster than a single group could.

\paragraph{Stage 6: deploy and self-correct.} Where affordable, methods are evaluated closer to deployment and made to correct themselves under the shift they meet there: by test-time training in non-reactive simulation (Chapter~\ref{ch:centaur}), by aligning training with execution on real robots (Chapter~\ref{ch:kai0}), or, in co-authored work, by splitting computation between a slow, knowledgeable model and a fast, reactive one~\cite{hamdan2025eta}.

\paragraph{Closing the loop.} The residual failures found in later stages feed the next iteration. Several iterations are visible. OpenLane was extended by OpenLane-V2 from lane lines to road topology, whose ground truth in turn generated the traffic-element questions of DriveLM. OpenOcc and the occupancy challenges led to OpenScene, which supplies the data of the NAVSIM benchmark on which Chapter~\ref{ch:centaur} is evaluated. DriveLM-Data and its challenge were audited by DriveBench~\cite{xie2025vlms}, which found that 78.6\% of DriveLM behaviour answers are ``Going Straight'' and that models can score without images, motivating more balanced data and more grounded metrics. In manipulation, the data-scaling platform AgiBot World was followed by \chiz, which studied alignment during post-training under resource constraints.

\paragraph{Coverage of the cycle.} Table~\ref{tab:method:stages} records which stages each project actually reached. No single project completed the whole cycle: the perception and reasoning projects stop at the community stage and are evaluated open-loop, whereas the acting projects reach deployment-like evaluation but have not (yet) been opened to a public challenge. The cycle is therefore best read as a description of a research programme spread over several projects, not of each project.

\begin{table}[t]
  \centering
  \caption[Stages of the research cycle reached by each project]{Stages of the research cycle (Fig.~\ref{fig:method:cycle}) reached by each project: (1)~observe, (2)~formalise, (3)~benchmark, (4)~baseline, (5)~community, (6)~deploy and self-correct. \checkmark: reached; $\circ$: partly reached; --: not reached. Projects marked $^{*}$ are co-authored rather than led by the author.}
  \label{tab:method:stages}
  \footnotesize
  \renewcommand{\arraystretch}{1.12}
  \begin{tabularx}{0.82\textwidth}{@{}X *{6}{>{\centering\arraybackslash}p{0.95cm}}@{}}
    \toprule
    Project & (1) & (2) & (3) & (4) & (5) & (6) \\
    \midrule
    PersFormer / OpenLane & \checkmark & \checkmark & \checkmark & \checkmark & -- & -- \\
    OpenLane-V2$^{*}$ & \checkmark & \checkmark & \checkmark & \checkmark & \checkmark & -- \\
    OccNet / OpenOcc (Ch.~\ref{ch:occnet}) & \checkmark & \checkmark & \checkmark & \checkmark & \checkmark & -- \\
    OpenScene (Ch.~\ref{ch:occnet}) & \checkmark & -- & \checkmark & $\circ$ & \checkmark & $\circ$ \\
    DriveLM (Ch.~\ref{ch:drivelm}) & \checkmark & \checkmark & \checkmark & \checkmark & \checkmark & -- \\
    DriveBench$^{*}$ & \checkmark & \checkmark & \checkmark & -- & $\circ$ & -- \\
    Centaur (Ch.~\ref{ch:centaur}) & \checkmark & \checkmark & \checkmark & \checkmark & $\circ$ & $\circ$ \\
    AgiBot World$^{*}$ & \checkmark & -- & \checkmark & \checkmark & \checkmark & \checkmark \\
    \chiz (Ch.~\ref{ch:kai0}) & \checkmark & \checkmark & $\circ$ & \checkmark & -- & \checkmark \\
    \bottomrule
  \end{tabularx}
\end{table}

\paragraph{Evidence and limits.} Two kinds of evidence suggest that the benchmark-driven part of the cycle has had an effect beyond the author's own projects: the uptake of its artefacts, such as the 152 teams of the ``Driving with Language'' track and the reuse of OpenScene as the data of NAVSIM, whose own CVPR 2024 competition drew 143 teams~\cite{dauner2024navsim}; and the cycle's ability to correct itself, as when DriveBench exposed shortcuts in DriveLM-Data. This evidence is observational: there is no counterfactual in which the same problems were studied without public benchmarks. The cycle is also not the only productive way to do research on physical agents. Hypothesis-driven work that tests a theoretical prediction, method-first work on well-established benchmarks, and deployment-first engineering in industry each have their place. The cycle is least appropriate when an adequate benchmark already exists, when the cost of building a benchmark exceeds its likely use, or when the failure cannot be captured in recorded data at all, as for events so rare that they must be generated or simulated rather than collected.

% ==============================================================================
\section{Research Lineage}\label{sec:method:lineage}

Table~\ref{tab:method:papers} summarises the eighteen publications to which the author contributed, and Fig.~\ref{fig:method:lineage} shows how they build on one another. Three observations stand out. First, the work moves along the perceive--reason--act loop over time, from 3D perception (2022--2023) to language-based reasoning (2023--2025) and to acting under shift (2025--2026). Second, each move is prepared by an artefact from the previous stage: the BEV representation of BEVFormer underlies both OccNet and UniAD; OpenScene supplies the data on which Centaur is evaluated; and UniAD is the driving baseline of DriveLM, whose driving-direction categories in turn anchor the uncertainty measure of Centaur. Third, benchmarks (dashed outlines) appear in every pillar rather than only at the start: each new pillar began with a new way of measuring. The figure mixes works that the author led with co-authored ones. Project leadership does not imply sole authorship of every component; the core chapters state the individual and shared contributions.

\begin{figure}[t]
  \centering
  \tikzfig{ch-methodology/figures/lineage}
  \caption[Research lineage of the author's publications]{Research lineage of the publications to which the author contributed, arranged by year and by pillar of the perceive--reason--act loop. Solid arrows indicate that a work builds on another (architecture, data or evaluation protocol); dashed arrows indicate that a work is motivated by, or contrasts with, another. Thick outlines mark the works led by the author that form the core chapters; dashed outlines mark works that release data or a benchmark.}
  \label{fig:method:lineage}
\end{figure}

\begin{table}[p]
  \centering
  \caption[Publications and their role in the research cycle]{The publications to which the author contributed, their pillar, their role in the research cycle of Fig.~\ref{fig:method:cycle}, the aspect of the curse of reality they address, and the author's role. ``Led'' denotes first, co-first or co-project-lead authorship.}
  \label{tab:method:papers}
  \footnotesize
  \renewcommand{\arraystretch}{1.10}
  \begin{tabularx}{\textwidth}{@{}>{\raggedright\arraybackslash}p{2.95cm} >{\raggedright\arraybackslash}p{1.9cm} >{\raggedright\arraybackslash}p{1.9cm} >{\raggedright\arraybackslash}X >{\raggedright\arraybackslash}p{1.5cm}@{}}
    \toprule
    Work (venue) & Pillar & Role in cycle & Aspect of reality addressed & Author \\
    \midrule
    LSH-SMILE~\cite{sima2021lshsmile} (NeurIPS'21) & earlier & method & cost of simulating and learning physical dynamics & first \\
    PersFormer, OpenLane~\cite{chen2022persformer} (ECCV'22) & perceive & benchmark, method & non-flat roads; 3D lanes from a single camera & co-first \\
    BEVFormer~\cite{li2022bevformer} (ECCV'22) & perceive & baseline & multi-camera 3D perception under occlusion & co-author \\
    BEV survey~\cite{li2023bevdevils} (T-PAMI'23) & perceive & study, recipe & sensor configurations; practical recipe (Waymo 2022, 1st) & co-first \\
    UniAD~\cite{hu2023uniad} (CVPR'23) & architecture & baseline & error accumulation in modular stacks & co-author \\
    SDF~\cite{gao2023sdf} (IROS'23) & perceive & method & information loss in fusion; sensor failure & co-author \\
    OpenLane-V2~\cite{wang2023openlanev2} (NeurIPS'23) & perceive & benchmark & road topology beyond lane lines & co-author \\
    \textbf{OccNet, OpenOcc}~\cite{sima2023occnet} (ICCV'23) & perceive & benchmark, method & irregularly shaped obstacles & \textbf{led} (Ch.~\ref{ch:occnet}) \\
    \textbf{OpenScene}~\cite{sima2023openscene} (2023) & perceive & benchmark & scale and diversity of occupancy data & \textbf{led} (Ch.~\ref{ch:occnet}) \\
    VCD~\cite{huang2023vcd} (NeurIPS'23) & perceive & method & LiDAR available at training time only & co-author \\
    \textbf{DriveLM}~\cite{sima2024drivelm} (ECCV'24) & reason & benchmark, baseline & unseen objects and sensor set-ups; explainability & \textbf{led} (Ch.~\ref{ch:drivelm}) \\
    Hint-AD~\cite{ding2024hint} (CoRL'24) & reason & method, dataset & explanations aligned with the driving model & co-author \\
    DriveBench~\cite{xie2025vlms} (ICCV'25) & reason & study, benchmark & reliability of VLM answers and metrics & co-author \\
    ETA~\cite{hamdan2025eta} (ICCV'25) & reason & method & latency of large models in the loop & co-author \\
    \textbf{Centaur}~\cite{sima2024centaur} (2025) & act & method, stress split & rare safety-critical scenes & \textbf{led} (Ch.~\ref{ch:centaur}) \\
    AgiBot World~\cite{bu2025agibot_iros} (IROS'25) & act & platform, data & scale and quality of robot data & co-author \\
    Self-directed learning~\cite{chen2025robot} (2026) & architecture & perspective & improvement beyond imitation & co-author \\
    \textbf{$\chi_0$}~\cite{sima2026kai0} (2026) & act & method, system & inconsistent demonstrations; train--deploy mismatch & \textbf{led} (Ch.~\ref{ch:kai0}) \\
    \bottomrule
  \end{tabularx}
\end{table}

% ==============================================================================
\section{Design Principles}\label{sec:method:principles}

Five design principles recur across the methods in Table~\ref{tab:method:papers}. They are stated here together with the evidence that supports them and the evidence that qualifies them.

\paragraph{DP1: represent the world without closing it.} Representations should avoid unnecessary restrictions on the shapes they can express. The progression from 3D boxes to bird's-eye-view features~\cite{li2022bevformer}, from lane lines to topology graphs~\cite{wang2023openlanev2}, and from boxes to semantic occupancy~\cite{sima2023occnet} follows this principle. Occupancy also transfers across tasks: in the open-loop planning experiment of Chapter~\ref{ch:occnet}, planning on occupancy instead of 3D boxes lowers the collision rate by 15--22\% across horizons, although the absolute differences are small (below 0.3 percentage points). The principle is qualified in two ways: dense 3D representations are expensive to annotate, which is why the data effort of OpenScene was needed; and occupancy opens geometry but not semantics, since its labels still come from a fixed set of classes.

\paragraph{DP2: make intermediate reasoning inspectable.} When a decision may use knowledge absent from task labels, intermediate predictions can be exposed for supervision and inspection. DriveLM structures driving as perception, prediction and planning QAs~\cite{sima2024drivelm}, and the co-authored Hint-AD aligns explanations with intermediate model outputs~\cite{ding2024hint}. In Chapter~\ref{ch:drivelm}, most of the displayed zero-shot advantage over a driving-specific model comes from the pre-trained backbone; selected graph-context conditions have smaller gains under shift and none in-domain, and graph topology is not isolated. DriveBench further shows that language outputs can look competent without visual grounding~\cite{xie2025vlms}. Inspectability is therefore a design property, not proof of causal reasoning or generalisation.

\paragraph{DP3: optimise for the final task.} Every module should be judged by its contribution to acting well. The co-authored UniAD made this the organising principle of an end-to-end driving system~\cite{hu2023uniad}; its ablations showed that planning benefits only when motion forecasting and occupancy prediction are both present. The same principle appears in the planning experiments that validate occupancy (Chapter~\ref{ch:occnet}) and in the behaviour and motion questions that end every DriveLM graph (Chapter~\ref{ch:drivelm}).

\paragraph{DP4: close the loop at deployment.} Training data may leave deployment conditions uncovered, motivating mechanisms for correction. The NAVSIM comparisons in Section~\ref{sec:method:shifts} concern the particular safeguards evaluated there, not safety filters as a class. Chapter~\ref{ch:centaur} instead adapts the planner with one gradient step on an uncertainty proxy without querying new expert labels, raising the NAVSIM driving score from 90.3 to 92.6 in a single run (92.10 $\pm$ 0.33 over three seeds). Chapter~\ref{ch:kai0} closes the loop on the training side, by collecting recovery data with targeted recovery collection and aligning the timing of training and execution. The co-authored perspective of~\cite{chen2025robot} generalises this principle into self-directed learning, discussed in Chapter~\ref{ch:discussion}.

\paragraph{DP5: respect the resource budget.} Methods must fit the data, compute and latency budget of deployment. This principle appears as efficiency in simulation~\cite{sima2021lshsmile}, as a tenfold reduction of the nuPlan data volume in OpenScene~\cite{sima2023openscene}, as the fast and slow split of the co-authored ETA, which reaches a closed-loop driving score of 69.53 at \SI{50}{\milli\second} per step against 64.22 at \SI{931}{\milli\second} for the previous best~\cite{hamdan2025eta}, and in the post-training setting of \chiz (Chapter~\ref{ch:kai0}), which builds on a heavily pre-trained policy; its total data and compute accounting remains to be reconciled. It is qualified by the scaling results of AgiBot World, where more and better data did pay off~\cite{bu2025agibot_iros}; the two views are discussed together in Chapter~\ref{ch:discussion}.

% ==============================================================================
\section{Evaluation Methodology}\label{sec:method:evaluation}

\subsection{A ladder of evaluation fidelity}

Evaluating a physical agent involves a trade-off between fidelity to reality and cost. Four levels of fidelity are distinguished (Fig.~\ref{fig:method:ladder}). At level L0, models are evaluated open-loop on logged data with per-frame metrics such as mIoU, mAP, question-answering scores or trajectory displacement error. L0 is cheap and reproducible but blind to the policy's own mistakes; for planning in particular, open-loop metrics on nuScenes can be matched by models that exploit the ego vehicle's own state without perceiving the scene~\cite{zhai2023rethinking,li2024ego}. At level L1, non-reactive simulation scores a planned trajectory by rolling it out against logged agents and checking it against rules for collisions, drivable area, progress and comfort; NAVSIM~\cite{dauner2024navsim} is the benchmark used in Chapter~\ref{ch:centaur}. At level L2, closed-loop simulation lets other agents react and errors compound, for example in the CARLA simulator~\cite{dosovitskiy2017carla}. At level L3, the policy runs on physical hardware and is judged by task success and recovery, as in Chapter~\ref{ch:kai0}.

\begin{figure}[t]
  \centering
  \tikzfig[0.92\textwidth]{ch-methodology/figures/evaluation_ladder}
  \caption[A ladder of evaluation fidelity]{A ladder of evaluation fidelity. Higher levels expose the policy more to the consequences of its own actions, at a higher cost per evaluation. The right-hand column lists where each level is used in this thesis; closed-loop simulation (L2) was used only to validate the rule-based expert that generates DriveLM-CARLA.}
  \label{fig:method:ladder}
\end{figure}

The guiding rule is to evaluate each claim at the highest level of fidelity that is affordable for it, and to state the level explicitly. Perception claims (Chapter~\ref{ch:occnet}) are made at L0, complemented by an open-loop planning experiment that asks whether the representation helps downstream. Reasoning claims (Chapter~\ref{ch:drivelm}) are made at L0, on nuScenes, Waymo and logged CARLA data. Driving-policy claims (Chapter~\ref{ch:centaur}) are made at L1, and manipulation claims (Chapter~\ref{ch:kai0}) at L3. No learned driving agent in this thesis was evaluated at L2 or on a real vehicle.

\subsection{Testing for shift}

Because the subject of this thesis is robustness under shift, in-distribution accuracy is necessary but not sufficient. Chapters~\ref{ch:drivelm}--\ref{ch:kai0} include evaluations designed to expose shift: zero-shot transfer across sensor configurations (nuScenes to Waymo) and to an unseen object class in Chapter~\ref{ch:drivelm}; a safety-critical split drawn from the periphery of the test distribution in Chapter~\ref{ch:centaur}; and operation from varied initial states on real robots in Chapter~\ref{ch:kai0}. Chapter~\ref{ch:occnet} evaluates irregularly shaped objects of known classes and the downstream use of occupancy, but does not test distribution shift; this limitation is discussed in Chapter~\ref{ch:discussion}.

\subsection{Metric validity and statistical practice}

Metrics are themselves artefacts that can fail. DriveBench showed that language-based driving metrics can reward answer templates rather than grounded reasoning~\cite{xie2025vlms}, which is why DriveLM reports behaviour and motion metrics alongside language metrics, and why the CVPR 2024 challenge combined a GPT-based score with language, matching and accuracy terms. Statistical practice in the core chapters follows the conventions of the field rather than best practice: most results are single training runs, as is common for large models, and no formal significance tests were performed. Where variability is available it is reported, for instance the three-seed mean of Centaur's leaderboard entry (92.10 $\pm$ 0.33), and real-robot results are obtained over repeated trials. Small differences should be read with this in mind, and the chapters flag the cases where a margin is within the likely run-to-run variation.

% ==============================================================================
\section{From a Failure to a Contribution}\label{sec:method:decision}

The framework and the cycle can be combined into a diagnostic guide for choosing which lever to investigate when a new failure is observed (Fig.~\ref{fig:method:decision}). The first question is whether the failure can be measured reproducibly; if not, the first contribution is a benchmark. The next questions locate the failure: whether the information needed to avoid it is missing from the representation (enrich the representation); whether the decision requires knowledge beyond the training labels (add structured reasoning); whether the situation is under-covered by the training data (scale or curate data); whether a usable label-free corrective signal is available at deployment (investigate test-time adaptation); and whether training and execution conditions differ (align training with deployment). Several levers may apply to one failure; the order is a reading aid, not a validated rule for choosing the best intervention. Whatever the lever, the result is evaluated at the highest affordable fidelity, released, and exposed to the community, and the residual failures start the procedure again. Reading the core chapters through this procedure explains their order: Chapter~\ref{ch:occnet} answers the representation question, Chapter~\ref{ch:drivelm} the reasoning question, and Chapters~\ref{ch:centaur} and~\ref{ch:kai0} the two deployment questions, while each begins by answering the measurement question.

\begin{figure}[p]
  \centering
  \tikzfig[0.86\textwidth]{ch-methodology/figures/lever_decision_revised}
  \caption[From an observed failure to a research contribution]{From an observed failure to a research contribution: a diagnostic guide distilled from the author's work. Each right-hand branch suggests a methodological lever; several may apply to the same failure. The order is not an execution sequence or a validated intervention-selection rule; the representation, reasoning, test-time and alignment levers are the subjects of Chapters~\ref{ch:occnet}--\ref{ch:kai0}, and the data lever is illustrated by the co-authored AgiBot World ($^{*}$). Every branch ends with evaluation, release and exposure to the community, whose residual failures restart the procedure.}
  \label{fig:method:decision}
\end{figure}

% ==============================================================================
\section{Reproducibility, Ethics and Data Governance}\label{sec:method:practice}

\subsection{Availability of artefacts}

Openness is the mechanism through which Commitment 1 operates, so it is treated as part of the method. Table~\ref{tab:method:artefacts} lists the public availability of the artefacts of the core chapters at the time of writing. The benchmarks are derived from established public datasets and simulators, including nuScenes~\cite{caesar2020nuscenes}, the Waymo Open Dataset~\cite{sun2020waymo}, nuPlan~\cite{caesar2021nuplan} and CARLA~\cite{dosovitskiy2017carla}, and are redistributed under the licences of their sources. Each chapter and its appendix report available hyper-parameters, hardware and training details; unresolved protocol and budget information is identified there.

\begin{table}[t]
  \centering
  \caption[Public availability of the artefacts of the core chapters]{Public availability of the artefacts of the core chapters at the time of writing.}
  \label{tab:method:artefacts}
  \footnotesize
  \renewcommand{\arraystretch}{1.15}
  \begin{tabularx}{\textwidth}{@{}p{1.5cm}p{4.1cm}X@{}}
    \toprule
    Chapter & Artefact & Availability \\
    \midrule
    \ref{ch:occnet} & OccNet code and OpenOcc data & \url{https://github.com/OpenDriveLab/OccNet} \\
    \ref{ch:occnet} & OpenScene data and development kit & \url{https://github.com/OpenDriveLab/OpenScene} \\
    \ref{ch:drivelm} & DriveLM-Data, baseline and evaluation server & \url{https://github.com/OpenDriveLab/DriveLM} \\
    \ref{ch:centaur} & Centaur code and the \texttt{navsafe} split & not publicly released at the time of writing \\
    \ref{ch:kai0} & \chiz code (Apache 2.0) & \url{https://github.com/OpenDriveLab/KAI0} \\
    \ref{ch:kai0} & \chiz data and checkpoints (CC BY-NC-SA 4.0) & Hugging Face repository \texttt{OpenDriveLab-org/Kai0} \\
    \bottomrule
  \end{tabularx}
\end{table}

\subsection{Ethics and data governance}\label{sec:method:ethics}

The studies reported here did not recruit people as experimental subjects, but people produced annotations and robot demonstrations. The available project records describe annotators verifying occupancy labels and writing or checking DriveLM-nuScenes QAs, and teleoperators collecting demonstrations and recovery data for \chiz. This thesis does not independently establish the institutional ethics classification, consent process, employment arrangements, privacy review or licence coverage for those activities; these must be checked against the original project records before submission. The underlying datasets have their own terms of use, including non-commercial restrictions for \chiz data, and derived releases must follow them. Some evaluation uses commercial language models whose versions may change and whose use incurs reproduction costs. The driving-data providers report privacy protections for nuScenes, nuPlan and Waymo. Centaur was evaluated only in simulation; real-road online adaptation would require separate safety and certification review.

% ==============================================================================
\section{Threats to Validity}\label{sec:method:validity}

Several threats qualify the conclusions of this thesis, and each is revisited where it applies.

\paragraph{Construct validity.} Benchmarks are proxies of reality, and metrics are proxies of what matters. Occupancy mIoU does not measure safety; the PDM score of NAVSIM is computed from non-reactive rollouts; language metrics may reward style over substance. The thesis mitigates this by pairing proxy metrics with downstream or closed-loop evaluations where possible, but cannot eliminate it.

\paragraph{Internal validity.} Automatically generated labels contain errors, and a benchmark designed by the authors of a method may favour that method. Human verification of labels, public release of benchmarks, and challenges in which other teams compete on the same data are the main safeguards.

\paragraph{External validity.} Results obtained in simulation or on specific datasets may not transfer to other vehicles, cities, robots or tasks. The zero-shot and safety-critical evaluations of Chapters~\ref{ch:drivelm} and~\ref{ch:centaur} probe transfer directly, and the real-robot evaluation of Chapter~\ref{ch:kai0} removes the simulation gap for manipulation, though only for the tasks and hardware studied.

\paragraph{Conclusion validity.} Leaderboard comparisons are subject to random variation and to the ``benchmark lottery'' of which tasks and datasets are chosen~\cite{dehghani2021benchmark}. Most results in the core chapters are single runs, so small margins may not be meaningful.

\paragraph{Hindsight.} The framework of this chapter was reconstructed after the research was done. Reconstructions tend to make a sequence of opportunistic decisions look more deliberate than it was, and the examples that fit a framework are easier to recall than those that do not. Table~\ref{tab:method:stages} records where the projects did not follow the cycle, as a partial guard against this bias.

% ==============================================================================
\section{Chapter Summary}\label{sec:method:summary}

This chapter has described the research framework behind the thesis. Perceptual, semantic, behavioural and deployment shift were organised as non-exclusive diagnostic categories, illustrated by methodological levers and core chapters, with data and measurement cutting across them. The research was described as a benchmark-driven cycle that begins with a failure in the field, turns it into a measurable task and a public benchmark, establishes a simple and scalable baseline, engages the community through open release and challenges, and ends with evaluation closer to deployment and self-correction, whose residual failures start the next iteration; no single project completed the whole cycle. Five design principles, a ladder of evaluation fidelity and a diagnostic guide complete the framework, which was reconstructed in hindsight and should be read with the threats of Section~\ref{sec:method:validity} in mind. The next four chapters present the core contributions, starting with the perception end of the loop: how a physical agent should represent the 3D world it sees.
